\chapter*{KẾT LUẬN}
\addcontentsline{toc}{chapter}{KẾT LUẬN}

Tiểu luận đã xây dựng một kế hoạch quản trị dự án tương đối đầy đủ cho đề tài ``Quản lí dự án gây dựng quỹ từ thiện bằng Blockchain'' trong bối cảnh môn học Quản trị dự án. Trọng tâm của báo cáo không nằm ở việc mô tả sâu công nghệ blockchain, mà ở cách tổ chức, kiểm soát và đánh giá một dự án phần mềm có nhiều ràng buộc về phạm vi, tiến độ, chi phí, chất lượng và rủi ro.

Từ việc phân tích bài toán minh bạch dòng tiền, báo cáo đã xác định được phạm vi sản phẩm và phạm vi dự án theo hướng khả thi trong ba tháng thực hiện. Các nội dung cốt lõi như Project Scope Statement, WBS, Schedule Baseline, Cost Baseline, kế hoạch chất lượng, kế hoạch nguồn lực, truyền thông, rủi ro, mua sắm và tích hợp đã được xây dựng theo cùng một logic xuyên suốt: mỗi deliverable phải có owner, tiêu chí chấp nhận, bằng chứng kiểm tra và cơ chế kiểm soát thay đổi.

Một điểm quan trọng của đề tài là sự phân biệt rõ giữa \textit{Product Scope} và \textit{Project Scope}. FundChain được xem là sản phẩm thử nghiệm nhằm minh họa quy trình gây quỹ minh bạch trên Sepolia, còn báo cáo này là hồ sơ quản trị giúp trả lời các câu hỏi: phải làm gì, làm trong bao lâu, tiêu tốn bao nhiêu nguồn lực, kiểm soát bằng cách nào và khi nào thì được xem là hoàn thành. Cách tiếp cận đó giúp tránh tình trạng đồng nhất ``có mã nguồn'' với ``dự án đã hoàn tất''.

Qua các chương, tiểu luận cũng cho thấy khi áp dụng quản trị dự án vào một sản phẩm có yếu tố blockchain, nhiều kiến thức quản trị truyền thống vẫn giữ nguyên giá trị nhưng cần diễn giải sát bối cảnh hơn. Ví dụ, gas trở thành biến số chi phí cần theo dõi; ABI, contract address và cấu hình môi trường trở thành configuration item; transaction hash, event log và CID trở thành bằng chứng nghiệm thu; còn rủi ro thay đổi smart contract sau triển khai buộc nhóm phải siết chặt change control và quality gate.

Về mặt học thuật, báo cáo giúp hệ thống hóa kiến thức môn học thành một trường hợp ứng dụng cụ thể, có đủ đầu vào để thảo luận về đánh đổi giữa phạm vi, thời gian, chi phí và chất lượng. Về mặt thực tiễn, bộ tài liệu này có thể được dùng làm nền tảng để tổ chức lại việc phát triển, kiểm thử, demo và bàn giao sản phẩm FundChain theo cách chuyên nghiệp và dễ truy vết hơn.

Bên cạnh các kết quả đạt được, tiểu luận vẫn có những giới hạn nhất định. Dự án được lập kế hoạch trên cơ sở thời gian môn học ba tháng và ngân sách giả định, do đó các chỉ số EVM, NPV hay cost reserve chủ yếu phục vụ mục tiêu học tập. Sản phẩm cũng mới dừng ở môi trường Sepolia, chưa trải qua audit bảo mật độc lập, đánh giá pháp lý, kiểm thử tải và kế hoạch vận hành production. Vì vậy, báo cáo không xem đây là một hệ thống sẵn sàng triển khai thực tế với tiền thật, mà là một case study quản trị dự án có tính ứng dụng.

Trong giai đoạn tiếp theo, nếu tiếp tục phát triển sản phẩm, nhóm nên ưu tiên bốn hướng: hoàn thiện hồ sơ phân vai cá nhân; tăng chiều sâu của kiểm thử và review bảo mật; bổ sung dữ liệu vận hành thực tế để hiệu chỉnh lịch biểu và chi phí; và đánh giá lại phạm vi nếu muốn tiến tới môi trường Mainnet hoặc có tích hợp pháp lý. Những bước này sẽ giúp chuyển dự án từ mức thử nghiệm học thuật sang mức sẵn sàng hơn cho ứng dụng thực tế.

Tổng thể, tiểu luận đã đạt mục tiêu cốt lõi của môn học: chứng minh rằng một dự án phần mềm dù mang màu sắc công nghệ mới vẫn phải được quản trị bằng các nguyên tắc nền tảng về lập kế hoạch, phân công trách nhiệm, kiểm soát chất lượng, quản trị rủi ro và kết thúc dự án một cách có hệ thống.
